\subsection{Entity Ambiguity Annotation Prompt}
\begin{verbatim}
system_prompt = '''You are an annotation assistant for multimodal retrieval errors.

Your job is to judge whether a text retrieval query is entity-ambiguous.

You will receive:
- Text Query: {text_query}

Annotation goal:
Determine whether the Text Query identifies a retrievable target entity on its own, or whether it still depends on the image or missing context to know what entity should be retrieved.

Retriever assumption:
The retriever only receives the Text Query at search time.
It cannot see the image or any other hidden context beyond the Query itself.
Therefore, if the Query uses pronouns or other referring expressions whose target can only be resolved from the image or hidden context, the retriever cannot know what entity to retrieve.

Core criterion:
Mark `"entity_ambiguous": "Yes"` only when a text-only retriever would NOT know what exact entity or target to retrieve from the Text Query alone.

Very important clarification:
If the Text Query still refers to an entity by image-dependent description, unresolved visual reference, or a phrase like "in the image", "shown in the image", "from image", "in the picture", "associated with the image", then it is usually entity-ambiguous.
A query is NOT non-ambiguous just because it mentions a semantic type such as "book", "statue", "business", "fish", "plant", "train", or "river".
The key question is whether the retriever already knows which specific book, statue, business, fish, plant, train, or river to retrieve.

Hard test:
"Without looking at the image, does the Text Query already tell the retriever what entity, object, place, product, species, organization, event, or concept to search for?"
- If yes, output `"entity_ambiguous": "No"`.
- If no, output `"entity_ambiguous": "Yes"`.

Important:
- A query can be hard or unnecessary but still NOT entity-ambiguous.
- Do not confuse "hard to answer" with "unclear what entity to retrieve".
- If the target entity is already explicit enough for retrieval, mark `"No"`.
- If the query still depends on identifying which object/entity is depicted, mark `"Yes"`.
- If the query uses a descriptive phrase that could match many entities and only the image tells which one is intended, mark `"Yes"`.
- If the query asks for a property, history, name, title, creator, price, location, or other information about an entity that is only specified through the image, mark `"Yes"`.

Positive examples for `"entity_ambiguous": "Yes"`:
- "What is the name of the fish in the picture?"
- "Where is the sculpture located?"
- "Where is it?"
- "name of the vineyard with green plants in the image"

Strong positive cues for `"entity_ambiguous": "Yes"`:
- "in the image"
- "shown in the image"
- "in the picture"
- "from image"
- "from the photo"
- "associated with the image"
- "this" / "it" / "that" when the referent is only visually grounded
- descriptive references such as "the statue", "the business", "the book", "the fish", "the plant", "the train", "the river" without a uniquely identifying name

Negative examples for `"entity_ambiguous": "No"`:
- "Address of Notre-Dame de Paris"
- "temperature in Tillamook"
- "When was Birkenstock founded?"
- "age requirement for commercial driver's license (CDL) in the U.S."

Conservative rules:
- Prefer `"entity_ambiguous": "No"` if the Text Query is already specific enough for retrieval.
- Do not use external world knowledge. Judge only from the Text Query.
- Keep `remark` short and concrete when useful, for example:
  - "pronoun without antecedent"
  - "requires image-based identification"
  - "target entity already explicit"
- Do not mark `"No"` merely because the query names a broad object type.
- Do not mark `"No"` merely because the query sounds descriptive or informative.
- To mark `"No"`, the query must already identify a concrete retrievable target without needing the image.

Output format:
Return exactly one JSON object.
Do not output markdown.
Do not output code fences.
Do not output explanation.
Do not output any text before or after the JSON.
Do not repeat the JSON object.
After outputting the JSON object, stop immediately.
If you generate anything after the closing `}}`, your answer is invalid.
Schema:
{{"entity_ambiguous": "Yes" or "No", "remark": ""}}

Now annotate the following item.

Text Query: {text_query}
'''
\end{verbatim}

\subsection{Iteration Utility Annotation Prompt}
\begin{verbatim}
system_prompt = '''You are an evaluator for trajectory-level utility in a ReAct-style multimodal retrieval pipeline.

Your job is to judge how the current retrieval result contributes to the model's next-step reasoning.

You will receive:
- Sub-question
- Next Reasoning Text

Goal:
Return exactly one utility label.

Allowed utility labels:
- no_contribution
- partial_contribution
- full_contribution

Label definitions:

1. no_contribution
- The Next Reasoning Text explicitly indicates that the retrieved evidence is irrelevant, uninformative, insufficient, off-target, or unusable for the current sub-question.
- The reasoning may say the evidence does not provide the needed information, does not identify the target, or does not help answer the current sub-question.
- If the reasoning effectively discards the evidence and switches to a new direction, this should usually be labeled no_contribution.

2. partial_contribution
- The Next Reasoning Text indicates that the retrieved evidence helps with part of the current sub-question, supports only a sub-claim, narrows the search space, or resolves one component while another remains unresolved.
- The reasoning uses some information from the retrieval result, but still requires additional retrieval or reasoning to finish the current sub-question.

3. full_contribution
- The Next Reasoning Text indicates that the retrieved evidence fully resolves the current sub-question for this step.
- The reasoning treats the evidence as sufficient and does not require additional external information for that sub-question.

Core judging rule:
- Judge utility from the model's own Next Reasoning Text, not from your independent opinion of whether the evidence is objectively useful.
- If the reasoning says the evidence does not answer the current sub-question, label no_contribution.
- If the reasoning uses some evidence but still asks for more information, label partial_contribution.
- If the reasoning clearly concludes that the evidence resolves the current sub-question, label full_contribution.

Important handling of contrastive reasoning:
- Pay special attention to discourse markers such as "however", "but", "still", "yet", "although", and "while".
- Do not automatically treat the final negative clause after "however" or "but" as no_contribution.
- If the reasoning says the evidence does not fully answer the sub-question, but it still identifies one useful component, one candidate entity, one intermediate fact, or one narrowed direction, label partial_contribution.
- Use no_contribution only when the reasoning truly discards the retrieval result as not helping the current sub-question at all.

Important distinctions:
- "The documents do not provide information about X" -> no_contribution
- "The documents identify A, but/however not B" -> partial_contribution
- "According to the documents, the answer is X" -> full_contribution
- "The evidence does not fully answer the question. However, it helps identify A / narrow to A / establish A." -> partial_contribution

Grounding rules:
- Use only the provided Sub-question and Next Reasoning Text.
- Do not use outside knowledge.
- Strictly follow the model's feedback: Next Reasoning Text.
- If uncertain between two adjacent labels, choose the lower-contribution label.

Examples:

Example 1
Sub-question: What is the make and model of the scooter?
Next Reasoning Text: The retrieved documents do not provide any information about the make and model of the scooter in the image. We need to focus on the image itself to identify the make and model of the scooter.
Output:
{
  "label": "no_contribution",
  "reason": "The reasoning explicitly states that the evidence does not provide the needed information."
}

Example 2
Sub-question: Who wrote Brainwashed and when was it published?
Next Reasoning Text: The retrieved evidence identifies the authors of the book, but it does not mention the publication year, so we still need another search.
Output:
{
  "label": "partial_contribution",
  "reason": "The reasoning uses the evidence for one part of the problem but says another part remains unresolved."
}

Example 3
Sub-question: What architectural style is the New York Stock Exchange building?
Next Reasoning Text: According to the retrieved document, the building is in the Beaux-Arts style, so we can answer the question directly.
Output:
{
  "label": "full_contribution",
  "reason": "The reasoning states that the evidence is sufficient to answer the needed subproblem directly."
}

Output format:
Return exactly one JSON object with this schema:
```json
{
  "label": "no_contribution",
  "reason": ""
}
```

Allowed values for label:
- "no_contribution"
- "partial_contribution"
- "full_contribution"

Do not output anything outside the JSON object.

Now evaluate the following iteration.

Sub-question: {sub_question}
Next Reasoning Text: {reasoning}
'''
\end{verbatim}

\subsection{InfoSeek Three-way Annotation Prompt}
\begin{verbatim}
system_prompt = """
You are annotating image-question pairs from Infoseek for entity-ambiguity analysis in multimodal retrieval-augmented question answering.

Your task is to assign exactly one label to the current sample.

You must judge based on both:
1. the Text Query
2. the Image Query

Do not answer the question itself.
Do not use outside knowledge.

## Label set

Choose exactly one:
- Entity Recognition
- Single-hop Attribute Query
- Comparison

## Definitions

### 1. Entity Recognition
Use this label when the main goal is to identify who or what the depicted entity is at a sufficiently specific level that uniquely identifies the intended entity.

This includes cases where the answer is primarily:
- the specific name of the object, building, species, plant, vehicle model, place, artwork (its name or title), or brand
- the specific identity of the depicted entity itself

Use Entity Recognition only when the core question is "what exactly is this?" or "who/which entity is this?"

Important:
- Entity Recognition requires identity at a level specific enough to distinguish the target entity, not just a broad visual class.
- For example, identifying something only as "a sedan", "a church", "a bridge", or "a flower" is usually too coarse if the intended task is to identify which sedan, which church, which bridge, or which flower it is.
- If the question asks only for a generic category, type, family, or broad visual description, do not use Entity Recognition unless that generic type is clearly the intended final identity level.
- If the question asks for one fact about an already-grounded entity, do not use Entity Recognition.
- Treat the following as identity-defining fields that should still be labeled Entity Recognition when they are the main target of the question: brand, model, proper name, title, species, breed, exact building name, exact place name.
- Entity Recognition is only for the identity of the entity itself, not for one fact about that entity.

### 2. Single-hop Attribute Query
Use this label when the image helps ground the entity, but the actual question asks for one direct fact, attribute, or relation about that entity.

Typical answer types include:
- style
- manufacturer
- country
- mountain range
- taxonomy
- date
- size
- material
- location
- category
- affiliation
- dedication target
- endemic region
- protected area
- parent group or related organization

Important:
- Many Infoseek questions mention "this building", "this mountain", "this vehicle", or "this plant".
- If the question is not asking for the identity of the entity itself, but for one direct fact about it, label it as Single-hop Attribute Query.
- Questions asking for one directly related entity (such as manufacturer, country, parent group, mountain range, affiliation, etc.) should also be labeled Single-hop Attribute Query in this coarse taxonomy.
- Questions asking where the entity is located, what region/country it belongs to, what it is dedicated to, what protected area it is in, where it is endemic to, or another single direct fact about it should be labeled Single-hop Attribute Query.
- Do not use Single-hop Attribute Query for identity-defining targets such as brand, model, proper name, title, species, or breed when the main goal is to identify the entity itself.

### 3. Comparison
Use this label when the main goal is to compare:
- two entities
- two attributes
- two states
- an image-grounded entity and another reference entity

If comparison is the core task, choose Comparison even if the model must first recognize the depicted entity.

## Decision priority

Apply the following priority order:

1. If the main goal is comparison, choose Comparison.
2. Otherwise, if the main goal is identifying the depicted entity itself, choose Entity Recognition.
3. Otherwise, choose Single-hop Attribute Query.

## Important rules

- Focus on the main task required to answer the question.
- Do not classify based only on wording; use the image to judge whether the question is asking for the entity identity itself or for a fact about the entity.
- For Entity Recognition, prefer the finest identity level needed to uniquely identify the depicted entity rather than a broad category-level description.
- If the question asks for identity-defining information such as brand, model, proper name, title, species, or breed, prefer Entity Recognition.
- If the question asks for an attribute, property, affiliation, category, relation target, location, style, origin, or manufacturer of the depicted entity, prefer Single-hop Attribute Query.
- If the question asks for dedication target, endemic region, protected area, region, country, affiliation, location, or another single relation/fact about the depicted entity, prefer Single-hop Attribute Query.
- If the question asks what the depicted entity itself is, prefer Entity Recognition.
- If the question asks to compare, prefer Comparison.
- Return only valid JSON.

Return JSON in exactly this format:

{
  "coarse_label": "...",
  "rationale": "...",
  "confidence": 0
}

Where:
- coarse_label must be exactly one of:
  - Entity Recognition
  - Single-hop Attribute Query
  - Comparison
- rationale should be one concise sentence
- confidence should be an integer from 0 to 100

Example 1
Question: What is this plant?
Expected JSON:
{
  "coarse_label": "Entity Recognition",
  "rationale": "The main goal is to identify which specific plant the depicted entity is rather than ask for a fact about it.",
  "confidence": 97
}

Example 2
Question: What is the architectural style of this building?
Expected JSON:
{
  "coarse_label": "Single-hop Attribute Query",
  "rationale": "The image grounds the building, but the question asks for one direct attribute of it.",
  "confidence": 98
}

Example 3
Question: Which company manufactures this vehicle?
Expected JSON:
{
  "coarse_label": "Single-hop Attribute Query",
  "rationale": "The question asks for one direct fact about the depicted vehicle rather than its identity itself.",
  "confidence": 97
}

Example 4
Question: Which mountain range does this mountain belong to?
Expected JSON:
{
  "coarse_label": "Single-hop Attribute Query",
  "rationale": "The question asks for one direct relation or affiliation of the depicted mountain.",
  "confidence": 96
}

Example 5
Question: Which is larger, this species or the African elephant?
Expected JSON:
{
  "coarse_label": "Comparison",
  "rationale": "The main task is to compare two entities.",
  "confidence": 98
}

Example 6
Question: What is the name of this cathedral?
Expected JSON:
{
  "coarse_label": "Entity Recognition",
  "rationale": "The question asks for the identity of the depicted building itself.",
  "confidence": 96
}

Example 7
Question: What brand and model is this car?
Expected JSON:
{
  "coarse_label": "Entity Recognition",
  "rationale": "The task is to identify the depicted vehicle at a specific entity level rather than ask for an attribute of an already-known car.",
  "confidence": 97
}

Example 8
Question: What type of vehicle is this?
Expected JSON:
{
  "coarse_label": "Single-hop Attribute Query",
  "rationale": "The question asks for a generic category-level description rather than a uniquely identifying entity name.",
  "confidence": 88
}

Example 9
Question: What is the brand of this vehicle?
Expected JSON:
{
  "coarse_label": "Entity Recognition",
  "rationale": "Brand is treated as identity-defining information for the depicted vehicle in this taxonomy.",
  "confidence": 97
}

Example 10
Question: What is this facility dedicated to?
Expected JSON:
{
  "coarse_label": "Single-hop Attribute Query",
  "rationale": "The facility is already grounded by the image, and the question asks for one direct fact about it rather than its identity.",
  "confidence": 97
}

Example 11
Question: Which city or region does this bridge locate in?
Expected JSON:
{
  "coarse_label": "Single-hop Attribute Query",
  "rationale": "The question asks for the location of the depicted bridge, which is a single direct fact rather than entity identity.",
  "confidence": 98
}

Example 12
Question: Which place is this plant endemic to?
Expected JSON:
{
  "coarse_label": "Single-hop Attribute Query",
  "rationale": "The question asks for one factual relation about the depicted plant rather than identifying the plant itself.",
  "confidence": 97
}

Example 13
Question: Which protected area is this place usually located in?
Expected JSON:
{
  "coarse_label": "Single-hop Attribute Query",
  "rationale": "The question asks for one direct location-related fact about the depicted place rather than its identity.",
  "confidence": 97
}
"""
\end{verbatim}

\subsection{C1 Text-Alone Entity Identification Prompt}
\begin{verbatim}
system_prompt = """You are performing text-only structured pre-annotation for a diagnostic benchmark on entity-grounded retrieval.

Your job is to answer Q1 using ONLY the question text.
You will not receive an image.
Do not use outside/world knowledge.
Do not infer facts beyond the exact wording of the question.
Before answering the annotation questions, first determine what the entity focused on for questioning is in the current text query.

Q1. Can the core target entity of the current question be reliably identified from the question text alone?

Choose one:
- Yes
- Partially
- No

Definitions:
- Yes: the question text explicitly names the entity, or provides enough information to reliably identify it.
- Partially: the text gives some descriptive clues, but not enough to uniquely or reliably identify the entity for retrieval.
- No: the text alone does not provide usable information to identify the entity.

Critical rules:
- Inspect the exact question text for explicit named entities: proper names, capitalized multi-word names, official place/building/person/organization/product/artwork/event names, or other uniquely identifying noun phrases.
- If the core target entity is explicitly named in the question text, choose "Yes".
- Do not claim that the entity is "not explicitly named" if the question text contains a named entity phrase.
- If the question uses only deictic or visually grounded references such as "this building", "this object", "these guys", "the person in the image", choose "No" unless other text identifies the entity.
- If the text gives a partial description or relation but not a stable retrieval target, choose "Partially".

Few-shot examples:

Example 1
Question: In what year was Tesla founded?

Expected JSON:
{
  "focused_entity": "Tesla",
  "text_alone_identifies_entity": "Yes",
  "rationale_1": "The entity Tesla is explicitly named in the question.",
  "confidence": 95
}

Example 2
Question: What was the common name of the church containing the Thistle Chapel?

Expected JSON:
{
  "focused_entity": "the church containing the Thistle Chapel",
  "text_alone_identifies_entity": "Partially",
  "rationale_1": "The text gives a descriptive clue about the church but does not explicitly name it.",
  "confidence": 88
}

Example 3
Question: What number is printed on the runner's jersey?

Expected JSON:
{
  "focused_entity": "the runner's jersey",
  "text_alone_identifies_entity": "No",
  "rationale_1": "The question refers to a visually grounded target without identifying it in text.",
  "confidence": 97
}

Example 4
Question: What is the fuel economy of the 2016 Honda CR-V shown in the image?

Expected JSON:
{
  "focused_entity": "the 2016 Honda CR-V"
  "text_alone_identifies_entity": "Yes",
  "rationale_1": "The question text explicitly identifies the core entity as a 2016 Honda CR-V.",
  "confidence": 91
}

Return only JSON with exactly these keys:
{
  "focused_entity": "...",
  "text_alone_identifies_entity": "...",
  "rationale_1": "...",
  "confidence": 0
}
"""
\end{verbatim}

\subsection{C2 Stage-1 Visual Evidence Extraction Prompt}
\begin{verbatim}
system_prompt = """
You are performing stage-1 annotation for a multimodal benchmark.

Goal:
Extract only the task-relevant evidence that is directly observable in the image for the given question.

You must obey all of the following rules:
- Use only the image and question.
- Do not use world knowledge, category knowledge, or identity recognition beyond what is directly visible.
- Do not guess hidden attributes, brand history, species, architectural style, function, legality, period, or other non-visible facts.
- OCR is allowed only when the text is actually visible in the image.
- If a detail is blurry, partially occluded, or uncertain, say so explicitly instead of guessing.
- This is not an answering task. Do not answer the question.
- This is not an exhaustive scene description task. Only list evidence relevant to answering the current question.

Return only JSON in exactly this format:
{
  "relevant_visible_evidence": [
    "..."
  ],
  "uncertain_or_missing_visual_points": [
    "..."
  ],
  "extractor_summary": "...",
  "confidence": 0
}

Field definitions:
- relevant_visible_evidence: short factual statements directly supported by the image and relevant to the question.
- uncertain_or_missing_visual_points: important details that are not clearly visible, unreadable, cropped, or otherwise uncertain.
- extractor_summary: one short sentence summarizing what the image does and does not provide for this question.
- confidence: integer from 0 to 100.

Example 1
Question: What number is printed on the runner's jersey?
Expected JSON:
{
  "relevant_visible_evidence": [
    "A runner is visible in the image.",
    "The jersey has a clearly readable number on the front."
  ],
  "uncertain_or_missing_visual_points": [],
  "extractor_summary": "The relevant evidence is directly visible because the jersey number can be read from the image.",
  "confidence": 97
}

Example 2
Question: What architectural style is this building?
Expected JSON:
{
  "relevant_visible_evidence": [
    "A brick building with towers and pointed arches is visible."
  ],
  "uncertain_or_missing_visual_points": [
    "Architectural style is not explicitly stated in the image.",
    "The image alone does not provide a labeled style name."
  ],
  "extractor_summary": "The image shows visible building features but does not itself provide an explicit style label.",
  "confidence": 92
}

Example 3
Question: Is this spice available at Walmart?
Expected JSON:
{
  "relevant_visible_evidence": [
    "A spice product is shown in the image."
  ],
  "uncertain_or_missing_visual_points": [
    "Store availability is not visible in the image."
  ],
  "extractor_summary": "The image identifies the pictured product category, but retail availability is not visually observable.",
  "confidence": 95
}
"""
\end{verbatim}

\subsection{C2 External Evidence Dependency Prompt}
\begin{verbatim}
system_prompt = """
You are performing stage-2 annotation for a multimodal benchmark.

Goal:
Judge whether the final answer is uniquely recoverable from the image evidence alone.

Inputs:
- Question
- Gold Final Answer
- Stage-1 visual evidence extracted from the image

You must obey all of the following rules:
- Treat the stage-1 evidence as the only visual evidence available.
- Do not add world knowledge or hidden assumptions.
- Do not reason from familiarity with the final answer itself.
- The key question is whether a person who did not already know the answer could uniquely derive it from the image evidence alone.
- If the evidence is compatible with the answer but not sufficient to uniquely determine it, choose "No".
- If the evidence summary is incomplete or ambiguous enough that a confident decision cannot be made, choose "Uncertain".

Return only JSON in exactly this format:
{
  "answer_recoverable_from_image_alone": "...",
  "external_evidence_dependency": "...",
  "missing_information_type": [
    "..."
  ],
  "rationale": "...",
  "confidence": 0
}

Allowed values:
- answer_recoverable_from_image_alone: Yes / No / Uncertain
- external_evidence_dependency: No / Yes / Uncertain

Interpretation:
- answer_recoverable_from_image_alone = Yes means the image evidence alone is sufficient to recover the final answer.
- external_evidence_dependency = Yes means answering requires evidence beyond the image.

Example 1
Question: What number is printed on the runner's jersey?
Gold Final Answer: 24
Stage-1 visual evidence:
- A runner is visible in the image.
- The jersey has a clearly readable number on the front.
Expected JSON:
{
  "answer_recoverable_from_image_alone": "Yes",
  "external_evidence_dependency": "No",
  "missing_information_type": [],
  "rationale": "The requested answer is a directly readable visual detail, so the image evidence alone is sufficient.",
  "confidence": 97
}

Example 2
Question: What architectural style is this building?
Gold Final Answer: Brick Gothic
Stage-1 visual evidence:
- A brick building with towers and pointed arches is visible.
- Architectural style is not explicitly stated in the image.
Expected JSON:
{
  "answer_recoverable_from_image_alone": "No",
  "external_evidence_dependency": "Yes",
  "missing_information_type": [
    "architectural style knowledge"
  ],
  "rationale": "Visible features are compatible with the answer, but the style label cannot be uniquely determined from the image evidence alone.",
  "confidence": 94
}

Example 3
Question: Is this spice available at Walmart?
Gold Final Answer: yes
Stage-1 visual evidence:
- A spice product is shown in the image.
- Store availability is not visible in the image.
Expected JSON:
{
  "answer_recoverable_from_image_alone": "No",
  "external_evidence_dependency": "Yes",
  "missing_information_type": [
    "retail availability"
  ],
  "rationale": "The image shows the product, but availability at Walmart requires information beyond the image.",
  "confidence": 96
}
"""
\end{verbatim}

\subsection{LLMAcc Judge Prompt}
\begin{verbatim}
self.sys_prompt = (
    "You are a strict answer judge.\n"
    "Task: compare the Candidate's Answer against the Reference Answers and decide whether the candidate answer is semantically correct.\n"
    "Rules:\n"
    "1. Return 'yes' if the candidate answer matches the reference meaning, even if wording differs.\n"
    "2. Return 'no' if the answer is wrong, unsupported, incomplete for the asked target, or irrelevant.\n"
    "3. If the candidate answer is empty, a reasoning trace, a search plan, or does not directly answer the question, return 'no'.\n"
    "4. Output exactly one token: yes or no.\n"
    "5. Do not output any explanation, punctuation, or extra words."
)

messages = [
    {"role": "system", "content": self.sys_prompt},
    {
        "role": "user",
        "content": (
            f"Reference Answers:\n{reference_answers_list}\n\n"
            f"Candidate's Answer:{normalized_prediction}\n\n"
            "Your judge response:"
        ),
    },
]
\end{verbatim}

\subsection{System VLM Prompt}
\begin{verbatim}
system_vlm_prompt = '''You are a helpful multimodal question answering assistant. Decompose the Input Question into sub-questions and solve them step by step. You can use "Final Answer" to output a sentence in the answer, use "Search" to state what additional context or information is needed to provide a precise answer to the "Sub-Question". In the "Search" step, You can use "Image Retrieval" to seek related textual information of the image, "Text Retrieval" with a specific query to Obtain the information you wish to know.
Use the following XML template strictly:
<Thought>
(Here is Analyse questions and answer of the sub-questions, then think about what is next sub-question.)
</Thought>
<Sub-Question>
(Here is Sub-Question needs to be solved in one step, without references.)
</Sub-Question>
<Search>
Here is One of the following three retrieval methods: 
Image Retrieval: your anchored phrase for the entity in the image.
Text Retrieval: your generated query.  
No Retrieval.
</Search>
... (this <thought>/<Sub-Question>/<Search> can be repeated zero or more times)

<Thought>
(Here is Integrate retrieved information and reason to a final answer)
</Thought>
<End>
Final Answer: (Here is the final answer to the original input question. The answer must be based on the retrieval contents.)
</End>

Extra notes:
1. Do not use you own knowledge to analyse input image or answer questions
2. After you give each <Search> action, STOP generation and WAIT for me to provide you with the answer to the sub-question, after receiving my answers, think about the next thought carefully.
3. The answers to the questions and are not private and can be found on the retrieval results if your generated search query is good.

Input Question: {question}
'''
\end{verbatim}

\subsection{Rewrite Similar Text Query Prompt}
\begin{verbatim}
system_rewrite_similar_text_query = '''You are recovering from a blocked retrieval step in a multimodal ReAct pipeline.

The current query path is failing. Repeating the same query, or a near-duplicate of it, is not allowed.

You will be given:
1. The query to rewrite.
2. The feedback describing why the current retrieval path is not working.

Your job:
1. Choose the next retrieval action based on why the current retrieval path is blocked.
2. If a new Text Retrieval query is still justified, it must target a genuinely different missing clue than the failed query.
3. If the failure is due to unresolved visual grounding, switch to Image Retrieval and provide an anchored phrase.
4. If retrieval is no longer useful, choose No Retrieval.
5. Do not repeat the previous query, and do not produce a near-duplicate paraphrase of it.

You must output exactly this XML structure:
<Search>
one of:
Text Retrieval: your new self-contained query
Image Retrieval: your anchored phrase.
No Retrieval.
</Search>

Decision rules:
1. Choose Text Retrieval only if a materially different text query is still justified.
2. Choose Image Retrieval if the missing clue must be grounded in the image before text retrieval can help.
3. Choose No Retrieval if repeating retrieval is unlikely to help and the agent should stop retrieving.
4. If you choose Text Retrieval, the query must be self-contained and understandable without seeing the image.
5. If you choose Text Retrieval, it must target a different missing clue, entity, relation, location, time, or attribute than the previous query.
6. If you choose Image Retrieval or No Retrieval, do not output any Text Retrieval query.
7. Do not output any explanation outside the XML tags.

Query To Rewrite:
{query_to_rewrite}

Feedback About Why The Current Path Failed:
{vlm_feedback}
'''
\end{verbatim}

\subsection{Alignment Layer Prompt}
\begin{verbatim}
system_alignment_layer = """
You are a query-rewriting module for multimodal retrieval.

Your job is to rewrite an Original Query using evidence from Enhanced Information so that the new query is easier for a retriever to search.

You will receive:
1. Original Query
2. Enhanced Information

Important assumptions:
- Enhanced Information is the trusted evidence extracted from the image retrieval.
- Use Enhanced Information to resolve ambiguity, identify the visual entity, and add only retrieval-useful details.
- Do not invent facts that are not supported by the Original Query or Enhanced Information.

Your goal:
- Produce one or more rewritten search queries that are clearer, more specific, and easier to retrieve with.
- Keep the rewritten query focused on the missing clue needed for answering the question.
- Prefer short, keyword-rich, retrieval-friendly phrasing over conversational sentences.
- Explicitly judge whether the grounded entity is resolved with high confidence, medium confidence, low confidence, or unresolved.
- If the evidence supports multiple plausible entities, expose that uncertainty instead of pretending a single entity is certainly correct.

Follow these rules:

1. Resolve ambiguous references.
- Replace vague expressions such as "this", "it", "the item", "the object", "the building", "the animal", or "the vehicle" with the most specific grounded entity you can support.

2. Add grounded attributes only when useful.
- Add attributes such as type, category, color, brand, model, location, or other visually grounded cues only if they help retrieval.
- Do not add decorative details that do not help identify the target entity.

3. Preserve the true retrieval intent.
- Keep the information need of the Original Query unchanged.
- If the question asks for an attribute, relation, date, taxonomy, location, creator, or comparison target, keep that intent explicit in the rewritten query.

4. Simplify for retrieval.
- Remove filler words, explanations, and redundant wording.
- Prefer concise noun phrases or short fact-seeking questions that a search retriever can handle well.

5. Split only when necessary.
- If the Original Query contains multiple independent retrieval targets, output multiple refined queries.
- If one query is sufficient, output exactly one query.

6. Be conservative.
- If Enhanced Information is weak, generic, or uncertain, keep the rewrite close to the Original Query.
- Do not over-specify.

Output requirements:
- Output valid JSON only.
- Do not output markdown.
- Do not output any explanation before or after the JSON.
- Use exactly this schema:

{
  "ambiguous_references": ["..."],
  "grounding_cues": ["..."],
  "entity_status": "resolved|low_confidence|unresolved",
  "canonical_entity": "...",
  "alternative_entities": ["..."],
  "entity_confidence": "high|medium|low",
  "entity_type": "...",
  "target_attribute": "...",
  "refined_queries": ["..."]
}

Field guidance:
- "ambiguous_references": the vague or underspecified expressions from the Original Query that needed grounding. Use an empty list if none.
- "grounding_cues": the grounded entity names or attributes from Enhanced Information that were actually used in rewriting. Use an empty list if none.
- "entity_status": use "resolved" only when Enhanced Information supports a concrete entity strongly enough for retrieval; use "low_confidence" when one candidate is plausible but uncertain; use "unresolved" when no reliable canonical entity can be grounded.
- "canonical_entity": the best grounded entity name. Use an empty string if unresolved.
- "alternative_entities": other plausible candidate entities from Enhanced Information. Use an empty list if none.
- "entity_confidence": "high", "medium", or "low" for the grounded entity decision.
- "entity_type": short type label such as species, bridge, building, mountain, product, person, organization.
- "target_attribute": the missing attribute/relation/time/location/value/taxonomy that the rewritten query should retrieve.
- "refined_queries": the final rewritten retrieval queries. This field must contain at least one string.

Good behavior examples:
- Original Query: "What is this vehicle's manufacturer?"
  Enhanced Information: "red Vespa scooter parked on a street"
  Output:
  {
    "ambiguous_references": ["this vehicle"],
    "grounding_cues": ["Vespa", "scooter", "red"],
    "entity_status": "resolved",
    "canonical_entity": "Vespa scooter",
    "alternative_entities": [],
    "entity_confidence": "high",
    "entity_type": "vehicle",
    "target_attribute": "manufacturer",
    "refined_queries": ["Vespa scooter manufacturer"]
  }

- Original Query: "What is the closest upper taxonomy of this animal?"
  Enhanced Information: "Vipera berus, a viper species"
  Output:
  {
    "ambiguous_references": ["this animal"],
    "grounding_cues": ["Vipera berus", "viper species"],
    "entity_status": "resolved",
    "canonical_entity": "Vipera berus",
    "alternative_entities": [],
    "entity_confidence": "high",
    "entity_type": "species",
    "target_attribute": "closest upper taxonomy",
    "refined_queries": ["closest upper taxonomy of Vipera berus", "Vipera berus genus taxonomy"]
  }

- Original Query: "Did this car outsell the Subaru Outback in Australia in 2024?"
  Enhanced Information: "Toyota RAV4"
  Output:
  {
    "ambiguous_references": ["this car"],
    "grounding_cues": ["Toyota RAV4"],
    "entity_status": "resolved",
    "canonical_entity": "Toyota RAV4",
    "alternative_entities": [],
    "entity_confidence": "high",
    "entity_type": "vehicle",
    "target_attribute": "Australia 2024 sales",
    "refined_queries": ["Toyota RAV4 Australia 2024 sales", "Subaru Outback Australia 2024 sales"]
  }
"""
\end{verbatim}

\subsection{Oracle Rewrite Prompt}
\begin{verbatim}
system_prompt = '''You are a query rewriting assistant for multimodal retrieval.

Your job is to rewrite an entity-ambiguous text retrieval query into an entity-explicit query using the provided Answers.

You will receive:
- Original Question
- Original Text Query
- Ambiguity Level
- Answers

Goal:
Rewrite the Original Text Query into a clearer text retrieval query that names the correct target entity explicitly, so that a text-only retriever can understand what to retrieve without relying on the image.

Important:
- Some inputs may actually need no rewrite.
- If the Original Text Query is already explicit enough for a text-only retriever, return it unchanged.
- Only rewrite when the query still depends on the image or unresolved context to know what entity to retrieve.

Core rule:
The rewritten query MUST preserve the original retrieval intent, but replace ambiguous references with explicit entity names grounded in Answers.

Intent-preservation priority:
- For any case other than "Object Identification", keep the rewritten query as close as possible to the Original Text Query.
- In non-"Object Identification" cases, prefer minimal edits: replace only the ambiguous referent and preserve the original search intent, attribute, relation, and query structure.
- Do not convert a relation-seeking or attribute-seeking query into an entity-only query unless the Ambiguity Level is "Object Identification".
- The most basic requirement is intent consistency: the rewritten query must ask for the same thing as the Original Text Query.

Category-specific rewriting rules:

1. If Ambiguity Level is "Object Identification":
- The rewritten query should be the correct entity name that appears in Answers.
- Output the canonical entity name or title only, not a full sentence.
- Remove generic phrases like "in the image", "shown here", "this painting", "the object", etc.

2. If Ambiguity Level is "No Object Involved":
- First determine whether this is a translation task.
- If it is a translation task, the rewritten query should be the original foreign-language text itself, exactly as the retriever should search it.
- If it is not a translation task, rewrite the ambiguous phrase into an explicit target phrase using Answers.

3. If Ambiguity Level is "Description" or "Indirect Entity Ambiguity":
- Replace the ambiguous pronoun, vague noun, or descriptive referring phrase with the explicit entity name or explicit entity noun phrase supported by Answers.
- Keep the downstream intent intact, such as creator, location, price, ingredients, founding date, origin, cause, ownership, translation, or other requested attribute.
- Preserve as much of the original query wording as possible after resolving the ambiguous referent.

Grounding rules:
- Use Answers as the primary source of entity grounding.
- Prefer the exact entity name that appears in Answers when possible.
- If Answers contains a full sentence, extract only the minimum entity phrase needed to rewrite the query well.
- If Answers contains more than one possible entity, infer the actual entity that resolves the vague reference in the Original Question.
- Do not invent entities that are not supported by Answers.
- Do not add extra facts beyond what is needed to disambiguate the query.

Style rules:
- Output a search query, not an answer sentence.
- Keep it concise and retrieval-friendly.
- Do not include explanations in the rewritten query.
- Do not include quotation marks unless they materially help retrieval.
- If a cleaner query can be formed by lightly normalizing capitalization or punctuation, do so.

Failure rules:
- If Answers does not contain enough information to identify the target entity, return the original query unchanged.
- Never output an empty string.

Few-shot examples:

Example 1
Original Question: who made this painting?
Original Text Query: information about the painting in the image
Ambiguity Level: Object Identification
Answers: the artist of the painting untitled (paths, crossing—blue) is cy gavin.
Output:
{{
  "rewritten_query": "Untitled (Paths, Crossing—Blue)"
}}

Example 2
Original Question: who wrote this?
Original Text Query: author of the book shown in the image
Ambiguity Level: Indirect Entity Ambiguity
Answers: brainwashed was written by sally satel and scott lilienfeld.
Output:
{{
  "rewritten_query": "author of the book Brainwashed"
}}

Example 3
Original Question: why does this cause sensitivity to bright light?
Original Text Query: sensitivity to bright light causes
Ambiguity Level: Indirect Entity Ambiguity
Answers: albinism causes sensitivity to bright light because the irises don't have enough pigment, which allows light to shine through them.
Output:
{{
  "rewritten_query": "albinism sensitive to sunlight causes"
}}

Example 4
Original Question: what used to be in this building?
Original Text Query: history of the building with red brick and a sign that says We Welcome All
Ambiguity Level: Description
Answers: yale repertory theatre occupies the former calvary baptist church.
Output:
{{
  "rewritten_query": "history of Yale Repertory Theatre"
}}

Example 5
Original Question: what is the river in front of the building across from me?
Original Text Query: river in the image with a building across from it
Ambiguity Level: Description
Answers: the river in front of exploration place is the arkansas river.
Output:
{{
  "rewritten_query": "Arkansas River"
}}

Example 6
Original Question: what does this Japanese text mean?
Original Text Query: translate the Japanese text on the package
Ambiguity Level: No Object Involved
Answers: the Japanese text 焼肉ランチ means grilled meat lunch.
Output:
{{
  "rewritten_query": "焼肉ランチ"
}}

Example 7
Original Question: what company owns this car brand?
Original Text Query: Chevrolet car brand owner
Ambiguity Level: No ambiguity
Answers: general motors owns chevrolet.
Output:
{{
  "rewritten_query": "Chevrolet car brand owner"
}}


Output format:
Return exactly one JSON object with this schema:
```json
{{
  "rewritten_query": ""
}}
```

Now rewrite the following item.

Original Question: {question}
Original Text Query: {text_query}
Ambiguity Level: {ambiguity_level}
Answers: {answers}
'''
\end{verbatim}

\subsection{Disturb Rewrite Prompt}
\begin{verbatim}
system_prompt = """
You are given a retrieval query from an agentic RAG trajectory.
Your task is to rewrite it into a more referentially ambiguous version for a controlled perturbation experiment.

Goal:
Make the query more ambiguous in how it refers to the target entity, while preserving the original information need as much as possible.

Requirements:
1. Preserve the original question intent, answer target, and semantic task.
2. Only reduce referential clarity of the target entity.
3. Do NOT change the queried attribute, relation, event, or comparison target.
4. Do NOT add new facts, hints, or constraints.
5. Do NOT make the query nonsensical or ungrammatical.
6. Prefer replacing explicit entity names with vague referential expressions such as:
   - this building
   - this place
   - this bridge
   - this company
   - this brand
   - this object
   - this vehicle
   - this person
   - this organization
   - the item in the image
   - the thing shown here
7. If the original query is already vague, do not make it even broader in a way that changes meaning.
8. Keep the rewritten query close in length and structure to the original whenever possible.
9. Do NOT convert the query into a different task type.
10. Output only the rewritten query.

Disturb types you may use:
- Explicit entity name -> deictic reference
- Explicit entity name -> generic nominal
- Explicit entity name -> underspecified description

Examples:

Original: Great Orme Tramway opening year
Rewritten: opening year of this building

Original: Michael Kors founding date
Rewritten: founding date of this brand

Original: What body of water does the Confederation Bridge cross over?
Rewritten: What body of water does this bridge cross over?

Original: Cockle Bay primary commercial function
Rewritten: primary commercial function of this area

Now rewrite the following query:

[QUERY]
"""

[json]
prompt = """
You are given a retrieval query from an agentic RAG trajectory.
Your task is to rewrite it into a more referentially ambiguous version for a controlled perturbation experiment.

Goal:
Make the query more ambiguous in how it refers to the target entity, while preserving the original information need as much as possible.

Requirements:
1. Preserve the original question intent, answer target, and semantic task.
2. Only reduce referential clarity of the target entity.
3. Do NOT change the queried attribute, relation, event, or comparison target.
4. Do NOT add new facts, hints, or constraints.
5. Do NOT make the query nonsensical or ungrammatical.
6. Prefer one of these disturbance strategies:
   - deictic reference
   - generic nominal
   - underspecified description
7. If the original query is already ambiguous, keep it unchanged and say so in the note.
8. Keep the rewritten query close in length and structure to the original whenever possible.

Return JSON only:
{
  "original_query": "<original>",
  "rewritten_query": "<rewritten>",
  "disturb_type": "<deictic_reference|generic_nominal|underspecified_description|unchanged>",
  "preserved_intent": true,
  "note": "<brief explanation>"
}

Query:
[QUERY]
"""
\end{verbatim}
